\chapter{Alphas intrajournaliers appris par apprentissage automatique}\label{ch:s1:intraday-machine-learned-alphas}

Un modèle à gradient boosté entraîné sur quatorze heures de carnet d'ordres simulé prédit les cinq prochaines secondes du prix milieu avec un R-deux
hors échantillon (out-of-sample) de 25~\%, les trente secondes suivantes avec 5~\%, et les deux minutes suivantes pas du tout. Utilisé de façon naïve (en franchissant la fourchette à chaque prévision d'un cinquième
de tick ou davantage), le modèle à cinq secondes perd 0.6 tick par action~: la fourchette qu'il paie est plus grande que les mouvements qu'il prédit. Utilisé de façon passive, en rejoignant la
file du côté de la prévision et en attendant d'être exécuté, les mêmes prévisions rapportent environ un tiers de tick par action. Un \omterm{def:s1:intraday-machine-learned-alphas:alpha}{alpha intrajournalier} est une prédiction et une
règle d'exécution à la fois~; aucune des deux ne vaut quoi que ce soit seule. La construction est \texttt{firm.intraml}.

\section{Cibles et horizons}

\begin{definition}[Alpha intrajournalier, horizon de prévision]\label{def:s1:intraday-machine-learned-alphas:alpha}
Un \emph{alpha intrajournalier}\index{alpha intrajournalier} est une prévision de la variation du prix d'un titre sur des durées allant de quelques secondes à quelques heures, construite à partir des données de marché (le carnet d'ordres, les
transactions, titres apparentés) et négociée au cours de la journée. Son \emph{horizon de prévision}\index{horizon de prévision} est l'intervalle sur lequel la variation de prix est prévue, ce qui
détermine la fréquence des négociations de l'alpha et le coût d'exécution qu'il peut supporter.
\end{definition}

La cible de ce chapitre est la variation du prix milieu sur les 5, 30 et 120 secondes suivantes, exprimée en ticks, échantillonnée toutes les secondes sur vingt sessions d'une heure de
\texttt{firm.tape} (livre~7, chapitre~2), dont la configuration par défaut comprend une fenêtre d'annonces en milieu de séance. Sa fourchette vaut environ un tick (1.07 en moyenne) et
l'écart type de la cible vaut 0.69 tick à cinq secondes, 2.19 à trente secondes et 4.47 à deux minutes. Kolm, Turiel et Westray, qui prévoient 115 actions du Nasdaq
à plusieurs horizons, ont constaté que les prévisions propres à chaque action avaient un horizon effectif d'environ deux variations de prix moyennes~: au-delà, le carnet dit
guère. La bande synthétique est cohérente avec ce constat dans sa forme.

\section{Variables issues du carnet et de la bande}

Dix variables, toutes calculées à partir des données antérieures à la seconde de la prévision (le \cref{lst:s1:intraday-machine-learned-alphas:features})~: le déséquilibre au meilleur niveau du carnet, la
fourchette, le déséquilibre du flux d'ordres (OFI) de Cont, Kukanov et Stoikov sur les 1, 5 et 30 dernières secondes (livre~7, chapitre~8), le volume signé des transactions sur 5 et 30
secondes (livre~7, chapitre~9), la variation propre du milieu sur 5 et 30 secondes, et la distance du microprix au milieu. Les variables de flux d'ordres, normalisées par la
profondeur du carnet, sont celles que la littérature juge stationnaires~: Kolm et ses coauteurs rapportent que les modèles entraînés sur le flux d'ordres surpassent la plupart des modèles entraînés sur
les états bruts du carnet. Les tests unitaires vérifient la causalité par construction~: tronquer la bande à un instant donné laisse inchangée chaque variable qui le précède.

\section{Modèles et leur validation}

Deux modèles par horizon~: une régression ridge sur les variables standardisées, et un ensemble à gradient boosté (méthode de Friedman) de soixante arbres de régression de profondeur deux
découpés sur les quantiles des variables, écrit en quelques dizaines de lignes de NumPy. L'entraînement utilise les sessions 1 à 14 et le test les sessions 15 à 20~: des sessions entières
séparées les unes des autres, afin que les cibles qui se chevauchent (une cible à trente secondes à la seconde $t$ partage 29 secondes avec celle de $t + 1$) ne puissent pas fuir d'un côté à l'autre de la coupure, la purge
et l'embargo du livre~7, chapitre~20 étant appliqués à l'échelle des sessions.

\begin{center}
\small
\begin{tabular}{@{}p{6.6cm}ccc@{}}
\toprule
horizon & 5 secondes & 30 secondes & 120 secondes\\
\midrule
R-deux hors échantillon (out-of-sample)~: ridge~; arbres boostés & 21.1~\%~; 24.6~\% & 5.1~\%~; 5.1~\% & $-0.4~\%$~; $-1.5~\%$\\
coefficient d'information de rang (IC)~: ridge~; arbres boostés & 0.33~; 0.34 & 0.23~; 0.28 & 0.07~; 0.07\\
R-deux dans l'échantillon (in-sample), arbres boostés & 23.3~\% & 9.2~\% & 3.2~\%\\
\midrule
chaque prévision au-dessus de 0.2 tick, agressive (ticks par action) & $-0.60$ & $-0.48$ & $-0.87$\\
chaque prévision au-dessus de 0.2 tick, passive & 0.37 & 0.33 & $-0.02$\\
couplée~: agressive au-dessus de la fourchette, passive en dessous & 0.33 & 0.22 & $-0.24$\\
\bottomrule
\end{tabular}
\end{center}

Le R-deux à cinq secondes, d'un quart, est élevé au regard des marchés réels, et la raison en est la bande~: le livre~7 a montré que ses milieux tendent à suivre une tendance pendant quelques
secondes, propriété que le modèle de ce chapitre apprend. Ce qui se transpose, c'est la configuration d'un horizon à l'autre (une décroissance rapide de la prévisibilité, avec la prévision à trente secondes
déjà tombée à 5~\% et celle à deux minutes sans valeur hors échantillon alors qu'elle s'ajuste encore à 3.2~\% dans l'échantillon) ainsi que le faible gain des arbres sur la régression ridge
là où il existe un signal à trouver. Sirignano et Cont ont constaté, avec de l'apprentissage profond sur des milliards de cotations américaines, qu'un modèle entraîné sur l'ensemble des actions prédit mieux
que des modèles propres à chaque action, y compris pour des actions hors de son ensemble d'entraînement~: la formation du prix à partir du flux d'ordres est largement universelle. C'est un argument pour regrouper
les données de plusieurs instruments, ce que les vingt sessions de ce chapitre ne font qu'imiter.

\section{Coupler l'alpha à l'exécution}

\begin{definition}[Couplage avec l'exécution]\label{def:s1:intraday-machine-learned-alphas:coupling}
Le \emph{couplage avec l'exécution}\index{couplage avec l'exécution} est la conception des ordres d'une stratégie intrajournalière à partir de sa prévision~: franchir la fourchette seulement lorsque la
prévision dépasse le coût du franchissement, se placer passivement lorsqu'elle ne le dépasse pas, et dimensionner et annuler les ordres à mesure que la prévision évolue, de sorte que l'alpha et
l'exécution ne forment qu'une seule décision.
\end{definition}

La moitié inférieure du tableau constitue le propos du chapitre. Une prévision d'un cinquième de tick ne couvre pas une fourchette d'un tick~: franchir la fourchette à chacune de ces prévisions fait perdre 0.60
tick par action à cinq secondes. Les mêmes prévisions, utilisées pour choisir de quel côté du carnet rejoindre, rapportent 0.37 tick par action sur les ordres exécutés
(le \cref{lst:s1:intraday-machine-learned-alphas:passive}), parce qu'un ordre passif gagne la fourchette au lieu de la payer, et que la prévision oriente les exécutions
vers le côté favorable. Coupler les deux (ne franchir la fourchette que lorsque la prévision la dépasse, soit 324 fois sur six heures de test, et se placer sinon)
rapporte 0.33 tick par action à cinq secondes et 0.22 à trente secondes. À deux minutes, rien ne fonctionne~: la prévision n'a aucun pouvoir prédictif hors échantillon, et les deux styles d'exécution
perdent de l'argent. Les chiffres passifs sont optimistes sur un point (le modèle d'exécution, ou fill model, exécute un ordre au repos dès que le marché traverse son prix, sans tenir compte de sa
position dans la file~; le livre~7, chapitre~18 propose le rejeu tenant compte de la file) et honnêtes sur un autre (la sortie franchit la fourchette).

\begin{omfigure}
\begin{tikzpicture}
\begin{axis}[name=a,ybar,bar width=8pt,width=6.4cm,height=5.2cm,xlabel={\small horizon (secondes)},ylabel={\small R-deux hors échantillon (\%)},
  tick label style={font=\footnotesize},symbolic x coords={5,30,120},xtick=data,ymin=-3,ymax=28,grid=major,grid style={gray!25},
  enlarge x limits=0.3,legend style={font=\scriptsize,at={(0.97,0.97)},anchor=north east},table/col sep=comma]
\addplot[fill=black!25,draw=black!60] table[x=h,y=ridge]{figdata/strategies-1/14-intraday-machine-learned-alphas/r2.csv};
\addplot[fill=omDef!60,draw=omDef] table[x=h,y=gbm]{figdata/strategies-1/14-intraday-machine-learned-alphas/r2.csv};
\legend{ridge,arbres boostés}
\end{axis}
\begin{axis}[at={(a.east)},anchor=west,xshift=2.0cm,ybar,bar width=6pt,width=6.4cm,height=5.2cm,xlabel={\small horizon (secondes)},
  ylabel={\small P\&L par action (ticks)},tick label style={font=\footnotesize},symbolic x coords={5,30,120},xtick=data,ymin=-1,ymax=0.6,
  grid=major,grid style={gray!25},enlarge x limits=0.3,legend style={font=\scriptsize,at={(0.5,1.03)},anchor=south,legend columns=3},
  table/col sep=comma]
\addplot[fill=omThm!60,draw=omThm] table[x=h,y=naive_aggressive]{figdata/strategies-1/14-intraday-machine-learned-alphas/pnl.csv};
\addplot[fill=omDef!40,draw=omDef] table[x=h,y=passive]{figdata/strategies-1/14-intraday-machine-learned-alphas/pnl.csv};
\addplot[fill=omDef!80,draw=omDef] table[x=h,y=coupled]{figdata/strategies-1/14-intraday-machine-learned-alphas/pnl.csv};
\legend{agressif,passif,couplé}
\end{axis}
\end{tikzpicture}

\omcaption{Prévisions par apprentissage automatique du prix milieu de la bande synthétique. À gauche~: R-deux hors échantillon par horizon. À droite~: résultat par action des prévisions à arbres boostés
sur six heures de test, négociées de façon agressive, passive ou couplée (agressive seulement au-delà de la fourchette). Données~: \texttt{s1\_intraml}.}
\label{fig:s1:intraday-machine-learned-alphas:results}
\end{omfigure}

\section{Fiches de stratégie}

\begin{strategyfile}{Alpha de déséquilibre du carnet}\label{strat:s1:intraday-machine-learned-alphas:book}
\sfield{Qui paie, et pourquoi} Des traders dont les ordres révèlent une pression avant que les prix ne bougent~: le déséquilibre du carnet est la file qui sera consommée ensuite.
\sfield{Instruments et plateformes} Actions et contrats à terme liquides, dotés de carnets profonds et visibles.
\sfield{Signal} Déséquilibre au meilleur niveau et déséquilibre de profondeur, déséquilibre du flux d'ordres sur quelques secondes.
\sfield{Dimensionnement et exécution} Surtout passive, en se plaçant du côté favorable~; agressive seulement lorsque la prévision dépasse la fourchette.
\sfield{Coûts} Fourchette, frais et rabais~; position dans la file.
\sfield{Comment elle meurt} Des concurrents plus rapides qui lisent le même carnet~; une profondeur factice.
\sfield{Horizon, capacité, infrastructure} Quelques secondes~; faible capacité par titre~; systèmes co-implantés à faible latence.
\sfield{Backtest honnête} Un rejeu tenant compte de la file (livre~7, chapitre~18)~; la latence~; les frais.
\sfield{Sources} Cont, Kukanov et Stoikov (2014)~; la simulation de ce chapitre.
\end{strategyfile}

\begin{strategyfile}{Alpha de flux de transactions}\label{strat:s1:intraday-machine-learned-alphas:flow}
\sfield{Qui paie, et pourquoi} Les métaordres se fragmentent en de nombreuses transactions, dont les signes persistent (livre~7, chapitre~9).
\sfield{Instruments et plateformes} Comme pour le déséquilibre du carnet d'ordres.
\sfield{Signal} Volume signé des transactions sur quelques secondes à quelques minutes~; persistance des signes de la transaction.
\sfield{Dimensionnement et exécution} Comme pour le déséquilibre du carnet d'ordres.
\sfield{Coûts} Comme pour le déséquilibre du carnet d'ordres.
\sfield{Comment elle s'épuise} Des métaordres plus discrets~; des concurrents sur le même flux d'ordres.
\sfield{Horizon, capacité, infrastructure} De quelques secondes à quelques minutes~; le flux de transactions avec leur signe.
\sfield{Backtest en toute honnêteté} Signes de transaction tels qu'on aurait pu les connaître (Lee--Ready ou l'indicateur d'agresseur du flux), avec sa latence.
\sfield{Sources} Kolm, Turiel et Westray (2023) sur les entrées de flux d'ordres.
\end{strategyfile}

\begin{strategyfile}{Alpha d'avance multi-actifs}\label{strat:s1:intraday-machine-learned-alphas:lead}
\sfield{Qui paie, et pourquoi} Des instruments plus lents qui suivent des instruments plus rapides (une action après son contrat à terme sur indice, un \omterm{def:s1:index-rebalancing:fund}{fonds indiciel} coté (ETF) après ses titres constitutifs).
\sfield{Instruments et plateformes} Paires d'instruments liés, sur une même plateforme ou sur des plateformes différentes.
\sfield{Signal} Le mouvement récent du meneur, relatif à celui du suiveur (livre~7, chapitre~10).
\sfield{Dimensionnement et exécution} Négocier le suiveur dans le sens du meneur avant qu'il ne rattrape son retard.
\sfield{Coûts} La vitesse~: le gain n'existe que pour celui qui arrive le premier.
\sfield{Comment elle s'épuise} Concurrence sur la latence~; l'avance se réduit jusqu'à quelques microsecondes.
\sfield{Horizon, capacité, infrastructure} De la milliseconde à la seconde~; réseaux et colocalisation (les livres~10 et~11).
\sfield{Backtest en toute honnêteté} Horodatages issus de la même horloge~; latence réaliste entre plateformes.
\sfield{Sources} Sirignano et Cont (2019) sur des variables universelles, mises en commun~; aucun chiffre de performance n'a été vérifié.
\end{strategyfile}

\begin{strategyfile}{Exécution guidée par l'alpha}\label{strat:s1:intraday-machine-learned-alphas:execution}
\sfield{Qui paie, et pourquoi} La fourchette, perçue au lieu d'être payée, lorsque la prévision choisit où placer l'ordre passif.
\sfield{Instruments et plateformes} Tout instrument que l'entreprise négocie pour d'autres raisons, et les portefeuilles de teneur de marché.
\sfield{Signal} La prévision à court horizon, mise en regard de la fourchette et de la file.
\sfield{Dimensionnement et exécution} Ordres agressifs au-dessus de la fourchette, passifs en dessous~; annulation lorsque la prévision s'inverse.
\sfield{Coûts} Sélection adverse sur les exécutions passives~; limites d'annulation.
\sfield{Comment elle s'épuise} Elle ne s'épuise pas~; elle réduit le coût de négociation de chaque autre stratégie.
\sfield{Horizon, capacité, infrastructure} Quelques secondes~; un moteur d'exécution qui prend les prévisions en entrée.
\sfield{Backtest en toute honnêteté} Exécutions tenant compte de la file~; les markouts des exécutions passives (livre~7, chapitre~23).
\sfield{Sources} La simulation de ce chapitre (0.33 tick par action couplé à $-0.60$ en exécution agressive à cinq secondes).
\end{strategyfile}

\section{Tutoriel~: un pour cent, c'est beaucoup}\label{tut:s1:intraday-machine-learned-alphas:tut}

\begin{tutorial}
\textbf{Objectif.} Construire les variables à partir du tape synthétique, entraîner et valider deux modèles à trois horizons, puis négocier leurs prévisions de trois façons. \textbf{État final~:}
le tableau et la \cref{fig:s1:intraday-machine-learned-alphas:results}.
\begin{tutsteps}
\item \textbf{Variables}~: à chaque seconde, à partir des données jusqu'à cette seconde.
\omcode{code/firm/intraml/firm_intraml.py}{34}{59}{Variables de carnet d'ordres et de flux de transactions.}\label{lst:s1:intraday-machine-learned-alphas:features}
\item \textbf{Exécution passive}~: se placer au meilleur cours du côté de la prévision~; exécution seulement si le marché a traversé le cours~; sortie de l'autre côté de la fourchette.
\omcode{code/firm/intraml/firm_intraml.py}{154}{172}{Entrées passives avec un modèle d'exécution au contact.}\label{lst:s1:intraday-machine-learned-alphas:passive}
\item \textbf{Lancer} \texttt{scores(h)}, \texttt{naive(h)} et \texttt{trading(h)} pour 5, 30 et 120 secondes, et \texttt{fig\_intraml.py}.
\end{tutsteps}
\textbf{Pour aller plus loin.} Remplacer le modèle d'exécution au contact par le rejeu du livre~7 tenant compte de la file~; ajouter le second instrument de \texttt{simulate\_pair} comme variable
d'avance~; entraîner un seul modèle sur tous les horizons, l'horizon étant une variable.
\end{tutorial}

\section{Construction~: apprentissage automatique intrajournalier}\label{bld:s1:intraday-machine-learned-alphas:intraml}

\begin{build}
\bfield{Rôle} Variables causales intrajournalières, cibles par horizon, deux modèles avec validation au niveau des séances, et règles d'exécution couplées aux prévisions.
\bfield{Interface} \texttt{features(tape, step)}, \texttt{targets(tape, times, horizons)}, \texttt{ridge} et \texttt{ridge\_predict}, \texttt{gbm} et
\texttt{gbm\_predict}, \texttt{r2(y, yhat)}, \texttt{aggressive(pred, fwd, spread, threshold)}, \texttt{passive(tape, times, pred, h, threshold, wait)}.
\bfield{Règles} Variables calculées à partir des données jusqu'à chaque instant~; apprentissage et test séparés par des séances entières~; coûts pris en compte dans chaque P\&L.
\bfield{Tests d'acceptation} \texttt{code/firm/intraml/tests/}~: variables inchangées lorsque le futur est retiré~; les deux modèles retrouvent un signal injecté~; le
P\&L agressif, vérifié à la main.
\bfield{Pour aller plus loin} Un modèle d'exécution passive tenant compte de la file~; des variables multi-actifs~; des arbres plus profonds et un arrêt anticipé sur une séance de validation.
\end{build}

\begin{omsources}
\item R.~Cont, A.~Kukanov et S.~Stoikov, «~The price impact of order book events~», \emph{Journal of Financial Econometrics} 12(1), 2014.
\item P.~N.~Kolm, J.~Turiel et N.~Westray, «~Deep order flow imbalance~», \emph{Mathematical Finance} 33(4), 2023.
\item J.~Sirignano et R.~Cont, «~Universal features of price formation in financial markets~», \emph{Quantitative Finance} 19(9), 2019.
\item J.~H.~Friedman, «~Greedy function approximation: a gradient boosting machine~», \emph{Annals of Statistics} 29(5), 2001.
\end{omsources}

\section{Exercices}

\begin{exercise}[$\star$]\label{exo:s1:intraday-machine-learned-alphas:1}
Une prévision affiche un R-carré de 25~\% face à une cible dont l'écart type vaut 0.69 ticks. Quel est l'écart type de la prévision~?
\end{exercise}

\begin{exercise}[$\star$]\label{exo:s1:intraday-machine-learned-alphas:2}
Pourquoi une cible à trente secondes à la seconde $t$ ne peut-elle pas être validée sur un échantillon qui contient la seconde $t + 1$~?
\end{exercise}

\begin{exercise}[$\star$]\label{exo:s1:intraday-machine-learned-alphas:3}
Un aller-retour agressif coûte une fourchette, soit un tick. Quelle doit être la prévision, exprimée en ticks, pour atteindre le seuil de rentabilité~?
\end{exercise}

\begin{exercise}[$\star\star$]\label{exo:s1:intraday-machine-learned-alphas:4}
Pourquoi les exécutions passives rapportent-elles de l'argent avec des prévisions trop faibles pour payer la fourchette~? Qu'est-ce qui biaise à la hausse les résultats passifs du chapitre~?
\end{exercise}

\begin{exercise}[$\star\star$]\label{exo:s1:intraday-machine-learned-alphas:5}
Le modèle à deux minutes affiche 3.2~\% dans l'échantillon (in-sample) et $-1.5~\%$ hors échantillon (out-of-sample). Qu'est-il arrivé~?
\end{exercise}

\begin{exercise}[$\star\star$]\label{exo:s1:intraday-machine-learned-alphas:6}
Pourquoi normaliser les variables de flux d'ordres par la profondeur du carnet~?
\end{exercise}

\begin{exercise}[$\star\star\star$]\label{exo:s1:intraday-machine-learned-alphas:7}
\emph{Programmation.} Entraîner les modèles après avoir retiré des variables les propres variations passées du prix milieu, puis comparer le R-carré à cinq secondes. Que nous apprend la
différence sur le tape synthétique~?
\end{exercise}

\begin{exercise}[$\star\star\star$]\label{exo:s1:intraday-machine-learned-alphas:8}
\emph{Trouver la faille.} «~Notre modèle prédit la variation de prix de la seconde suivante avec 60~\% de justesse sur la direction~; nous le négocierons avec des ordres au marché sur 3\,000
actions.~»
\end{exercise}

\section{Problème~: un pour cent, c'est beaucoup}

\begin{problem}[{Problème du week-end --- une prévision et son exécution}]\label{pb:s1:intraday-machine-learned-alphas:1}
Les modèles du chapitre sur \texttt{firm.tape}, ainsi que l'historique publié.

\medskip\noindent\textbf{Partie I --- La prévision.}
\begin{enumerate}
\item Définir un \omterm{def:s1:intraday-machine-learned-alphas:alpha}{alpha intrajournalier} et son \omterm{def:s1:intraday-machine-learned-alphas:alpha}{horizon de prévision}.
\item Décrire les cibles, les séances et la fourchette.
\item Lister les variables et expliquer comment la causalité est testée.
\item Que Kolm, Turiel et Westray ont-ils trouvé au sujet des horizons et des entrées~?
\end{enumerate}

\medskip\noindent\textbf{Partie II --- Les modèles.}
\begin{enumerate}[resume]
\item Décrire les deux modèles.
\item Comment l'entraînement et le test sont-ils séparés, et pourquoi par séance~?
\item Donner le R² hors échantillon et l'IC par horizon.
\item Pourquoi le R² à cinq secondes est-il si élevé ici~?
\end{enumerate}

\medskip\noindent\textbf{Partie III --- Exécution.}
\begin{enumerate}[resume]
\item Définir le \omterm{def:s1:intraday-machine-learned-alphas:coupling}{couplage avec l'exécution}.
\item Donner le P\&L par action des trois styles d'exécution par horizon.
\item Pourquoi les transactions agressives perdent-elles sur les prévisions faibles~?
\item Que suppose le modèle d'exécution passif~?
\end{enumerate}

\medskip\noindent\textbf{Partie IV --- Le verdict.}
\begin{enumerate}[resume]
\item Énoncer le \emph{résultat nommé}~: R² hors échantillon par horizon et P\&L par action après la fourchette, en agressif contre passif.
\item Que Sirignano et Cont ont-ils trouvé au sujet de l'universalité~?
\item Comment mettriez-vous en commun les données entre instruments~?
\item Que changerait un rejeu tenant compte de la file~?
\item Comment la latence intervient-elle~?
\item Quelle fiche de stratégie est la moins exposée à la concurrence~?
\item Quelle est la capacité d'un alpha à cinq secondes~?
\item En une phrase~: que vaut un \omterm{def:s1:intraday-machine-learned-alphas:alpha}{alpha intrajournalier}~?
\end{enumerate}
\end{problem}

\section{Questions d'entretien}

\begin{interviewq}[$\star$ \iqroles{chercheur, ingénieur ML}]\label{iq:s1:intraday-machine-learned-alphas:1}
Quelles variables construiriez-vous à partir d'un carnet d'ordres pour prédire les quelques secondes suivantes~?
\end{interviewq}

\begin{interviewq}[$\star\star$ \iqroles{chercheur, ingénieur ML}]\label{iq:s1:intraday-machine-learned-alphas:2}
Comment validez-vous par validation croisée un modèle dont les cibles se chevauchent dans le temps~?
\end{interviewq}

\begin{interviewq}[$\star\star$ \iqroles{trader}]\label{iq:s1:intraday-machine-learned-alphas:3}
Votre modèle prédit des mouvements plus petits que la fourchette. Est-il inutile~?
\end{interviewq}

\begin{interviewq}[$\star\star$ \iqroles{ingénieur ML, développeur}]\label{iq:s1:intraday-machine-learned-alphas:4}
Comment serviriez-vous un modèle par gradient boosting en quelques microsecondes après une mise à jour des données de marché~?
\end{interviewq}

\begin{interviewq}[$\star\star$ \iqroles{chercheur}]\label{iq:s1:intraday-machine-learned-alphas:5}
Pourquoi un R² de 1~\% pourrait-il être un bon modèle intrajournalier~?
\end{interviewq}

\begin{interviewq}[$\star\star\star$ \iqroles{chercheur}]\label{iq:s1:intraday-machine-learned-alphas:6}
Une prévision $f$ d'un mouvement $y$ a une corrélation $\rho$ avec lui, et les deux sont normales d'écarts-types $\sigma_f$ et $\sigma_y$. Avec un coût aller-retour
$c$, établir le P\&L espéré par transaction en ne négociant que si $|f| > c$, et discuter sa dépendance en $\rho$.
\end{interviewq}
